%% ============================================================
%%  SECTION 6 — KẾT LUẬN (Conclusion)  (số liệu: v2)
%% ============================================================
\section{Kết luận}
\label{sec:conclusion}

\subsection{Tóm tắt quá trình và kết quả}

Đồ án đã hoàn thành trọn vẹn quy trình Khoa học Dữ liệu (KDD) 6 bước:
thu thập 133.586 bản ghi từ 10 nguồn, tích hợp về lược đồ thống nhất,
làm sạch và phân tích khám phá, xây dựng và đánh giá 5 mô hình hồi quy,
kết thúc bằng ứng dụng demo Streamlit.

Mô hình triển khai \textbf{XGBoost} (sau kiểm toán MNAR + ràng buộc đơn điệu,
Mục~\ref{subsec:audit}) đạt:
\begin{itemize}[leftmargin=*, itemsep=2pt]
  \item \textbf{Test R\textsuperscript{2} = 0.565}; \textbf{Test RMSE = 1.895 triệu/tháng};
        \textbf{Test MAE = 1.26 triệu/tháng} — kèm \textbf{bảo đảm cấu trúc}:
        tick thêm tiện ích không bao giờ làm giảm giá dự báo
        (0 hiệu ứng âm, đơn điệu tuyệt đối trên app).
  \item Để tham chiếu: Random Forest chính xác hơn (Test R\textsuperscript{2} = 0.664), và
        biến thể XGBoost chỉ-làm-sạch-dữ-liệu (không ràng buộc) còn cao hơn nữa (0.671) —
        nhưng cả hai đều mang hiệu ứng tiện ích âm phi lý. Khoảng cách $\approx$0.10
        R\textsuperscript{2} là cái giá \textit{có chủ đích} của việc từ chối tương quan
        phi nhân quả để đổi lấy hành vi đúng cho app (Mục~\ref{subsec:audit}).
  \item Kích thước $\approx$0.86 MB, triển khai được trong giới hạn GitHub và app Streamlit.
\end{itemize}

\subsection{Bài học kinh nghiệm và ý nghĩa}

\begin{enumerate}[leftmargin=*, itemsep=5pt]
  \item \textbf{Diện tích và khu vực là yếu tố định giá mạnh nhất.}
        Feature Importance của cả RF lẫn XGBoost đều xếp \texttt{area}
        và \texttt{city}/\texttt{district} dẫn đầu sau khi loại \texttt{source}.

  \item \textbf{Nghịch lý Simpson là bài học về phân tích dữ liệu nghiêm túc.}
        Tương quan âm toàn cục giữa tiện ích và giá không phản ánh quan hệ nhân quả —
        đây là confounding variable do cấu trúc nguồn.
        Phát hiện này dẫn tới quyết định loại \texttt{source}
        (đánh đổi một phần độ chính xác nhưng đổi lại mô hình \textbf{trung thực và tổng quát hơn}).

  \item \textbf{Chỉ số tốt chưa đủ — phải kiểm thử hành vi và kiểm toán nguồn gốc dữ liệu.}
        Một hiệu ứng phi lý trên app (tick máy lạnh làm giá \textit{giảm}) được truy vết
        về tận dữ liệu thô, phát hiện \textbf{7/12 cờ tiện ích bị điền ngẫu nhiên ở bước
        tích hợp} ($\approx$97\% số dòng) — điều mà mọi chỉ số tổng hợp (RMSE, R\textsuperscript{2},
        CV) đều không phát hiện được. Bài học kép: (i) dữ liệu hợp nhất đa nguồn cần
        \textit{lưu vết nguồn gốc} (provenance) và phân biệt ``thiếu'' với ``không có''
        (MNAR); (ii) mô hình triển khai cho người dùng cần \textbf{đúng về hành vi}
        (ràng buộc đơn điệu), kể cả khi phải trả giá bằng chỉ số trên tập Test
        cùng phân phối với dữ liệu lỗi.

  \item \textbf{Các quyết định tiền xử lý phải được kiểm chứng bằng thực nghiệm, không chỉ lập luận.}
        Ba ablation study xác nhận:
        \begin{itemize}[itemsep=2pt]
          \item FE (tách khu vực từ địa chỉ) giảm RMSE 3.5--9.4\%;
                PCA không giúp ích với bộ dữ liệu này.
          \item Log-transform mục tiêu \textit{không} tăng hiệu năng theo tiêu chí RMSE/R\textsuperscript{2}
                và gây underweight phân khúc giá cao.
          \item Lọc ngoại lệ giảm RMSE 83.1\% trên vùng giá mục tiêu.
        \end{itemize}

  \item \textbf{Quan hệ phi tuyến mạnh.}
        Mô hình cây (R\textsuperscript{2} $\approx$ 0.57--0.66) vượt xa tuyến tính
        (R\textsuperscript{2} $\approx$ 0.19), xác nhận tương tác
        khu vực $\times$ diện tích không thể biểu diễn bằng hàm tuyến tính đơn giản.

  \item \textbf{Giới hạn bởi bias (không phải variance).}
        Learning curve cho thấy thêm dữ liệu cùng loại giúp ít —
        cần đặc trưng \textit{mới} (GPS, khoảng cách, NLP mô tả) để tăng R\textsuperscript{2}.
\end{enumerate}

\subsection{Hướng phát triển (Future Work)}

\begin{itemize}[leftmargin=*, itemsep=3pt]
  \item \textbf{Thu thập lại tiện ích thật} cho các nguồn đang thiếu
        (\texttt{chothuenhadattoanvn}, \texttt{github\_*}) và chạy lại toàn bộ
        so sánh 5 mô hình trên pipeline đã sửa MNAR — gỡ giới hạn
        ``tương quan hành vi đăng tin'' tận gốc thay vì bằng ràng buộc.
  \item \textbf{Đặc trưng vị trí không gian:} tọa độ GPS, khoảng cách đến trung tâm,
        tiện ích công cộng (trường học, bệnh viện, metro) — các nghiên cứu quốc tế
        cho thấy đây là yếu tố tăng R\textsuperscript{2} mạnh nhất.
  \item \textbf{Khai thác mô tả dạng văn bản:} NLP (TF-IDF, BERT) trên trường
        \texttt{description} để trích xuất tín hiệu giá (tầng, hướng, view, trang trí…).
  \item \textbf{Cân bằng dữ liệu địa lý:} thu thập thêm mẫu ngoài Hà Nội
        để mô hình đại diện tốt hơn cho các tỉnh thành khác.
  \item \textbf{Cập nhật định kỳ:} triển khai pipeline tự động cào và retrain
        hàng tháng, theo dõi data drift khi thị trường thay đổi.
  \item \textbf{Mô hình phức tạp hơn:} LightGBM, CatBoost, hoặc deep learning
        (TabNet) nếu bổ sung được đặc trưng GPS và văn bản.
\end{itemize}
